\begin{figure}[htbp]
\centering
\begin{tikzpicture}[>=Stealth, line width=0.8pt, font=\small,
	nodo/.style={circle, fill=white, draw=black,
		inner sep=0pt, minimum size=5pt, line width=0.8pt},
	etasta/.style={circle, draw=black, fill=white,
		inner sep=1pt, font=\scriptsize, line width=0.6pt},
	richiam/.style={thin, gray, dashed}]
	
	\def\L{2.8}    % lunghezza campata
	\def\tw{0.22}  % semi-larghezza triangolo
	\def\thA{0.38} % altezza triangolo cerniera (A)
	\def\thR{0.32} % altezza triangolo carrello
	\def\rn{0.08}  % raggio semicerchio al vertice
	\def\bw{0.16}  % semi-distanza cerchietti carrello
	\def\gw{0.30}  % semi-larghezza barra di terreno
	\def\yd{-1.10} % quota linee di misura
	
	%% ── Trave ────────────────────────────────────────────────────
	\draw[line width=2.2pt] (0,0) -- ({3*\L},0);
	
	%% ── Etichette nodi ───────────────────────────────────────────
	\node[above=0.05] at (0,0)      {$A$};
	\node[above=0.05] at ({\L},0)   {$B$};
	\node[above=0.05] at ({2*\L},0) {$C$};
	\node[above=0.05] at ({3*\L},0) {$D$};
	
	%% ── Numerazione campate ───────────────────────────────────────
	\node[etasta] at ({0.5*\L},  0.38) {$1$};
	\node[etasta] at ({1.5*\L},  0.38) {$2$};
	\node[etasta] at ({2.5*\L},  0.38) {$3$};
	
	%%═══════════════════════════════════════════════════════════════
	%% CERNIERA in A
	%%═══════════════════════════════════════════════════════════════
	\draw[line width=0.8pt]
	(0,0)--++(-\tw,-\thA)--++(2*\tw,0)--cycle;
	\draw[line width=0.8pt]({-\gw},{-\thA})--({\gw},{-\thA});
	\foreach \x in {-0.24,-0.12,0,0.12,0.24}
	\draw[thin]({\x},{-\thA})--++(-.09,-.11);
	\node[nodo] at (0,0){};
	
	%%═══════════════════════════════════════════════════════════════
	%% APPOGGIO INTERMEDIO in B
	%%═══════════════════════════════════════════════════════════════
	\draw[line width=0.8pt]
	({\L},0)--++(-\tw,-\thR)--++(2*\tw,0)--cycle;
	\draw[line width=0.8pt]
	({\L-\gw},{-\thR})--({\L+\gw},{-\thR});
	\draw[line width=0.8pt]({\L-\bw},{-\thR-0.12}) circle(2.4pt);
	\draw[line width=0.8pt]({\L+\bw},{-\thR-0.12}) circle(2.4pt);
	\draw[line width=0.8pt]
	({\L-\gw},{-\thR-0.24})--({\L+\gw},{-\thR-0.24});
	\foreach \x in {-0.24,-0.12,0,0.12,0.24}
	\draw[thin]({\L+\x},{-\thR-0.24})--++(-.09,-.11);
	\draw[line width=0.8pt, fill=white]
	({\L-\rn},0) arc(180:360:\rn) -- cycle;
	
	%%═══════════════════════════════════════════════════════════════
	%% APPOGGIO INTERMEDIO in C
	%%═══════════════════════════════════════════════════════════════
	\draw[line width=0.8pt]
	({2*\L},0)--++(-\tw,-\thR)--++(2*\tw,0)--cycle;
	\draw[line width=0.8pt]
	({2*\L-\gw},{-\thR})--({2*\L+\gw},{-\thR});
	\draw[line width=0.8pt]({2*\L-\bw},{-\thR-0.12}) circle(2.4pt);
	\draw[line width=0.8pt]({2*\L+\bw},{-\thR-0.12}) circle(2.4pt);
	\draw[line width=0.8pt]
	({2*\L-\gw},{-\thR-0.24})--({2*\L+\gw},{-\thR-0.24});
	\foreach \x in {-0.24,-0.12,0,0.12,0.24}
	\draw[thin]({2*\L+\x},{-\thR-0.24})--++(-.09,-.11);
	\draw[line width=0.8pt, fill=white]
	({2*\L-\rn},0) arc(180:360:\rn) -- cycle;
	
	%%═══════════════════════════════════════════════════════════════
	%% CARRELLO in D
	%%═══════════════════════════════════════════════════════════════
	\draw[line width=0.8pt]
	({3*\L},0)--++(-\tw,-\thR)--++(2*\tw,0)--cycle;
	\draw[line width=0.8pt]
	({3*\L-\gw},{-\thR})--({3*\L+\gw},{-\thR});
	\draw[line width=0.8pt]({3*\L-\bw},{-\thR-0.12}) circle(2.4pt);
	\draw[line width=0.8pt]({3*\L+\bw},{-\thR-0.12}) circle(2.4pt);
	\draw[line width=0.8pt]
	({3*\L-\gw},{-\thR-0.24})--({3*\L+\gw},{-\thR-0.24});
	\foreach \x in {-0.24,-0.12,0,0.12,0.24}
	\draw[thin]({3*\L+\x},{-\thR-0.24})--++(-.09,-.11);
	\node[nodo] at ({3*\L},0){};
	
	%%═══════════════════════════════════════════════════════════════
	%% QUOTATURA
	%%═══════════════════════════════════════════════════════════════
	\draw[richiam](0,      {-\thA-0.15})  -- (0,      {\yd+0.08});
	\draw[richiam]({\L},   {-\thR-0.40})  -- ({\L},   {\yd+0.08});
	\draw[richiam]({2*\L}, {-\thR-0.40})  -- ({2*\L}, {\yd+0.08});
	\draw[richiam]({3*\L}, {-\thR-0.40})  -- ({3*\L}, {\yd+0.08});
	
	\draw[thin, gray, |<->|] (0,{\yd}) -- ({\L},{\yd});
	\node[above=0.04, font=\small] at ({0.5*\L},{\yd}) {$\ell_1$};
	\draw[thin, gray, |<->|] ({\L},{\yd}) -- ({2*\L},{\yd});
	\node[above=0.04, font=\small] at ({1.5*\L},{\yd}) {$\ell_2$};
	\draw[thin, gray, |<->|] ({2*\L},{\yd}) -- ({3*\L},{\yd});
	\node[above=0.04, font=\small] at ({2.5*\L},{\yd}) {$\ell_3$};
	
\end{tikzpicture}
\caption{Trave continua a tre campate di luci $\ell_1$, $\ell_2$,
	$\ell_3$ ($n_e = 3$, $n_a = 4$): cerniera in $A$,
	appoggi intermedi in $B$ e $C$, carrello in $D$.
	Il cerchio in $A$ e $D$ indica il perno della
	connessione; il semicerchio in $B$ e $C$ indica
	che la trave scorre sull'appoggio senza interruzione
	della continuità rotazionale.
	Gradi di libertà $n = 9$, vincoli $m = 11$,
	grado di iperstaticità $g_{\mathscr{i}} = 2$
	pari al numero di appoggi intermedi.}
\label{fig:trave_continua}
\end{figure}
